\documentclass[12pt]{article}
\usepackage{amssymb}
\usepackage{amsmath}
\usepackage{float}

\begin{document}
  
  You can type $\mathbb{R}$ by \textbackslash mathbb\{Character\}, and don't forget to \textbackslash usepackage\{amssymb\}, for example:\\
  The distributive states that $a(b+c)=ab+ac$ for all $a, b, c \in \mathbb{R}$.

  \bigskip
  It's easy to type squer brackets [R], but to type curly prackets you have to type backslash (\textbackslash)
  before each curly bracket, for example :\\
  The set $A$ is defined to be $\{1, 2, 3\}$

  \bigskip
  Insert the doller sign type $ \$ $.

  \bigskip
  $ 2(\dfrac{1}{x^2-1}) $ \hspace{0.5cm} Oh, the brackets is bad, let's fix them :
  $ 2\left(\dfrac{1}{x^2-1}\right) $
  $$ 2\left[\frac{1}{x^2-1}\right] $$
  $$ 2\left\{\frac{1}{x^2-1}\right\} $$
  $$ 2\left\langle \frac{1}{x^2-1} \right\rangle $$
  $$ 2\left| \frac{1}{x^2-1} \right| $$
  $ \dfrac{dx}{dy}|_{x=1} $ \hspace{5mm} Oh, the form of the $|$ is short, let's fix it:
  $ \left.\dfrac{dx}{dy}\right|_{x=1} $\\

\noindent\rule{15.5cm}{2pt}

\vspace{0.5cm}
  Tables:\\\\
  \begin{tabular}{|c||c|c|c|c|c|} % c for center aligned, l for left aligned, r for right aligned.
      % you can use any kind of combination like {rcclll}.
      \hline
      x & 1 & 2 & 3 & 4 & 5\\ \hline % we saperate colums with ampersand (&).
      $ f(x) $ & 10 & 20 & 30 & 40 & 50\\ \hline 
      
  \end{tabular}

\vspace{0.2cm}
\begin{table}[H] % The compiler put the table in the place where it fit best,
% and to tell the compiler to put the table where we typed it, we use [H] that requiers to \usepackage{float}.  
\centering
\def\arraystretch{1.5} % to increase the space form the top to the bottom of the table cells 
\begin{tabular}{|c||c|c|c|c|c|}
    \hline
    x & 1 & 2 & 3 & 4 & 5\\ \hline
    $ f(x) $ & $\frac{1}{2}$ & 20 & 30 & 40 & 50\\ \hline 
    
\end{tabular}
\caption{These values represent the function $f(x)$.}
% You can type the caption above or bellow the table. (here is bellow)
\end{table}


\begin{table}[H]
\centering
\caption{The relationship between $f$ and $f'$.}
\def\arraystretch{1.5} 
\begin{tabular}{|l|p{5cm}|} 
    \hline
    $f(x)$ & $f'(x)$\\ \hline
    $x>0$ & The function $f(x)$ is increasing.\newline
    I will type here a lot of words to make the sentence to long 
    to show you how to fit a paragraph in the cells of the column 
    now I think that's enough.\newline
    you can edite the width of the column from p\{here\}.\\
    \hline 
    
\end{tabular}
\end{table}

\vspace{1cm}
Arrays:
\begin{align}
  5x^2-9=x+3\\
  5x^2-x-12=0
\end{align}

\begin{align}
  5x^2-9&=x+3\\
  5x^2-x-12&=0
\end{align}

\begin{align*}
  5x^2-9&=x+3\\
  5x^2-x-12&=0
\end{align*}

\begin{align}
  5x^2 \text{place your words here}\\
  5x^2 \, \text{place your words here}\\
  5x^2 \hspace{2mm} \text{place your words here}
\end{align}

\end{document}